\documentclass[10pt]{article}
\usepackage[a4paper,margin=23mm]{geometry}
\usepackage{amsmath,amssymb,booktabs,graphicx,hyperref}
\setlength{\emergencystretch}{2em}
\title{Query-Conditioned Partial-Evidence Fusion for 3D Scene Queries}
\author{TMM working manuscript}
\date{October 16, 2026}
\begin{document}
\maketitle
\begin{abstract}
Foundation-model responses can identify an object without recovering its complete
extent. A geometrically aligned observation may refer to the wrong instance,
while an incomplete region cannot automatically identify its exterior as
background. We study query-conditioned partial-evidence fusion over frozen
representations: determine whether an observation applies to a query, which
locations it may support or exclude, and how its response should enter the
existing geometric readout. The retained implementation supplies positive
controlled evidence: query-conditioned admission improves LERF selected-set
projection IoU from 64.12 to 66.66, and using already admitted observations in
the semantic response improves IoU from 50.59 to 70.38. These are distinct local
controls, not additive gains of one new equation. Four complete development
studies across a 31-job cohort further distinguish evidence admission from
correction magnitude and reconstruction accuracy. Lower source objectives in
all 22 source-assisted jobs do not yield consistent task improvements. The
results support preserving effective identity and extent providers while
concentrating method development on evidence authority. Functional preservation
is verified; a common admission rule with cross-teacher and cross-query gains
remains the central open validation target.
\end{abstract}

\section{Problem and contribution}
A region containing a cup handle is useful positive evidence, but may omit the
rest of the cup. A larger region may include the cup and its supporting table.
Another region can be geometrically consistent yet correspond to a different
cup. These situations require different decisions about the evidence. Treating
all regions as complete masks, or simply fitting them more accurately, need not
produce a better response to the user's query.

We study one problem: \emph{which observations are entitled to modify an existing
3D query response, and over what support?} The mechanism is query-conditioned
partial-evidence fusion. Applicability, spatial authority and geometric fusion
are successive parts of that mechanism. Teacher representation, feature
compression, geometry learning and complete-member prediction are held fixed
so their known capabilities are retained.

The study is organized around three contributions. First, it formulates
response evidence with separate query applicability, positive support,
exclusion authority and unknown support. Second, it connects this formulation
to a working candidate-admission rule and function-preserving extraction of
existing response stages. Third, it provides controlled evidence separating
admission, update size and source reconstruction quality. The first two
historical positive controls support the mechanism locally; the complete
system reference establishes available capability. Neither substitutes for
matched cross-provider validation of the common admission rule.

\input{sections/method.tex}

\section{Evaluation contract and retained capability}
The complete cohort contains 31 jobs: four LERF scenes shared by semantic and
selected-set projection evaluation, eight NVOS requests, ten ScanNet rooms and
nine UQIS rooms. UQIS includes 67 targets and four input types, totaling 268
queries. Table~\ref{tab:system} reports one complete retained recipe, not
metric-wise or task-wise winners. Its table builder verifies all job identities,
recomputes macro metrics and binds the results to source hashes.

\begin{table}[ht]
\centering\small
\resizebox{\linewidth}{!}{\input{tables/system_reference.tex}}
\caption{Complete retained v10 capability reference. These results belong to
its actual providers and readouts, not to an unvalidated common refinement rule.}
\label{tab:system}
\end{table}

LERF, NVOS and ScanNet use declared source-assisted evidence. NVOS retains a
positive-derived reference-mask protocol; it is not equated with a different
signed-scribble protocol. UQIS uses map-only query inputs and inherits a member
model trained on 22 separate labelled rooms and calibrated on three others.
The system is therefore not a label-free universal query model. LERF selected
sets are evaluated by 2D projection, not independent 3D ground truth. Its
projection IoU is 66.96; ScanNet 19/15/10-class mIoU is 61.01/62.53/72.49.
UQIS text/image/2D-point/3D-point AP is 26.60/59.31/75.79/75.33.

All four new studies seal the complete cohort's predictions before scoring.
They reuse opened development scenes; bootstrap intervals describe that cohort
and are not independent final-test guarantees. Post-refinement studies use
the retained complete response as an explicit input. Response-stage studies
instead preserve pre-fusion providers and intervene where the response is
originally generated. The independent selected-set heads stay fixed. Their
unchanged scores do not establish an improvement in instance selection.

\section{Evidence for the mechanism}
\paragraph{Candidate identity before spatial use.}
The retained text control fixes geometry, features and the set solver, and
changes the candidate admission rule. Conditioning native-score comparison on
candidates for the same query improves selected-set projection IoU from 64.12
to 66.66, with all four scenes improving. The rule avoids treating weak absolute
native response as sufficient reason to reject a detected part. Its one-versus-two
score models remain fallible; they do not certify cross-view instance identity.

\paragraph{Admitted evidence must reach the output.}
A separate control keeps the same admitted evidence, map and actual selected
sets while changing the continuous response output. Visibility fusion improves
semantic IoU from 50.59 to 70.38. The large effect establishes that useful
available evidence was missing from that output. It is not wholly attributable
to a novel fusion equation, and it cannot be added to the identity-control gain.

\begin{table}[ht]
\centering\small
\resizebox{\linewidth}{!}{\input{tables/mechanism_evidence.tex}}
\caption{Historical controlled effects tied to the component actually changed.
The first two rows support the mechanism locally. The final two justify
preserving effective identity and complete-extent providers.}
\label{tab:mechanism}
\end{table}

\paragraph{Preserving identity and complete extent.}
UQIS image local correspondence improves actual instance IoU from 30.74 to
34.88 with teachers, geometry and extent head held fixed. Reusing complete
membership then reaches 39.90; the latter paired scene interval is
$[2.40,7.74]$ percentage points, with seven of nine scenes improving.
These effects justify retaining the providers rather than replacing them with
a weaker initial similarity. Point membership gains from real observation
training and ScanNet gains from a second native teacher are also provider
evidence; their detailed controls and supervision are in the supplement.
They are not counted as net gains of the shared admission mechanism.

\section{Controlled diagnosis and method decisions}
\label{sec:stage-study}
Table~\ref{tab:controls} condenses all completed response controls. Each condition
uses a fixed rule across the full cohort. No row is assembled from the best
scene-specific or query-specific outputs. Detailed equations, all metrics and
per-query diagnosis remain in the accompanying diagnostic supplement.

\begin{table}[ht]
\centering\small
\resizebox{\linewidth}{!}{\input{tables/condensed_controls.tex}}
\caption{Complete development controls (\%). L: LERF; N: NVOS; S: ScanNet;
U: UQIS AP by query type. Post conditions act on the retained final response;
stage conditions act on frozen pre-fusion providers. $\dagger$: unchanged UQIS
compatibility output, not a reconstruction intervention. Independent selected
sets remain fixed throughout.}
\label{tab:controls}
\end{table}

\paragraph{Admission is not a one-sign update restriction.}
The first post-refinement study permits partial foreground support but excludes
unlicensed negative updates. Active-residual normalization nevertheless lowers
LERF IoU to 57.62 and degrades all UQIS modalities. Available-domain normalization
recovers LERF to 70.93, with a mean gain of 0.53 points and descriptive interval
$[0.18,0.89]$, while AP falls and other tasks still regress. Thus preserving
confidence and observable capacity matters, but the positive-only restriction
is insufficient. It is a diagnostic hypothesis, not the final method definition.

\paragraph{Extent agreement must outperform a smaller update.}
The third study weights admitted observations by agreement with a frozen extent.
Its matched control preserves each query's positive and negative pre-clipping
L1 update masses using uniform attenuation. Scope agreement is worse on UQIS
text/image/2D-point/3D-point AP by 0.50/0.87/0.59/0.88 points. It is not promoted
as a new module. Earlier evaluator-only diagnosis also finds 65 correctly
located supported anchors in each point modality, yet 48 and 46 of those
queries lose AP after refinement. Wrong anchor identity cannot be the sole
explanation; positive spatial support itself can be inaccurate. Labels in
this diagnosis never become query-time evidence permissions.

\paragraph{More accurate reconstruction is not sufficient.}
A source-lineage audit finds that all families already consume the same kind
of evidence upstream. The fourth study therefore moves the intervention into
response generation. Diagonal, rendered squared and rendered Huber solvers
share all admitted targets and weights. Both iterative arms use 120 steps and
step 0.95; Huber clips residuals at 0.25, inherited from NVOS. ScanNet retains
both native priors, and UQIS complete membership is unchanged.

Squared source objectives decrease in all 22 source-assisted jobs, with median
relative reductions of 19.64\%, 20.94\% and 1.50\% for LERF, NVOS and ScanNet.
Nevertheless LERF AP falls in all four scenes and ScanNet 19-class mIoU falls
in eight of ten rooms. Huber recovers 0.446 ScanNet points over squared loss,
improving all ten rooms, but remains 0.286 points below diagonal fusion.
This isolates an insufficient optimization target, not a proof that the loss
is caused entirely by identity, partial extent, geometry or view inconsistency.

\begin{figure}[ht]
\centering
\includegraphics[width=\linewidth]{figures/stage_source_fit_vs_task.pdf}
\caption{Per-job source objective reduction and task score change for rendered
squared reconstruction relative to diagonal consensus. Every source objective
improves, while downstream effects differ. Panels have separate horizontal
scales; the two NVOS horns requests are separate points.}
\label{fig:source-fit}
\end{figure}

\section{Implementation and remaining validation}
\paragraph{Estimated identity versus a forced positive.}
An additional causal audit removes only the positive override at each of the
67 estimated image anchors. All patches are sealed before scoring all 268 UQIS
queries. Raw classifier probabilities remain 0.914--0.939, above the frozen 0.35
threshold, so every actual set is unchanged. Image AP changes by only
0.00014 percentage points. Of 23 incorrect anchors, 16 lie in selected regions
that also contain target elements. This evaluator-only stratification motivates
testing identity transfer within mixed observations, but does not establish a
causal explanation for all residual error. Other benchmarks are retained
compatibility references in this audit, not freshly inferred transfer results.

The extracted fusion and reconstruction cores are connected to the original
production entry points. In the latest 22-source-job replay, original-solver
threshold/category predictions are unchanged, with maximum response deviations
$4.77\times10^{-7}$, $4.72\times10^{-5}$ and $7.75\times10^{-7}$ for LERF,
NVOS and ScanNet. NVOS maximum mean error is $6.78\times10^{-8}$; tiny AP drift
is retained as numerical variation. Targeted tests cover numerical
behavior, partial permissions and transient-run recovery. These checks establish
functional preservation, not the new admission rule's accuracy.

The next decisive experiment must change query-conditioned admission at the
actual response-generation stage while preserving identity, extent, input
access and the numerical backend. It must include an unconditioned control and
an evidence- or update-budget-matched control. UQIS must receive an actual
legal map-evidence intervention before it counts as transfer, and at least two
appropriate teacher configurations must be tested without changing the shared
rule per dataset. New validation scenes must remain untouched until the rule
is frozen. Improving an independent selected-set output requires a separately
measured intervention to that output.

The response study's dense predictions are volatile and released after scoring;
code snapshots, input/output hashes and complete metrics remain. The full study
occupies about 1.90 GiB under its 2 GiB budget. Storage policy and compatibility
engineering support reproducibility but are not presented as algorithmic
contributions. The retained complete recipe remains the performance reference;
none of the diagnostic alternatives has replaced it as a universal rule.

\input{sections/related_work.tex}
\end{document}
